\section{Theoretical Analysis}
\label{sec:theory}

We derive unified error bounds for the three-layer compression pipeline, connecting quantization error, eviction error, and information-theoretic limits.

\subsection{Quantization Error Bound}

\begin{theorem}[Quantization Error Bound]
\label{thm:quant}
For TurboQuant with $b_K$-bit Keys and $b_V$-bit Values applied after Hadamard rotation, the expected squared output error per attention head is bounded by:
\begin{equation}
    E_{\text{quant}} \leq d\sigma_V^2 \left( c_{b_V} \kappa + \frac{c_{b_K} \sigma_Q^2 \sigma_K^2}{T^2} \right)
\end{equation}
where $\kappa = \|a\|_2^2$ is the attention concentration, $T$ is the effective softmax temperature, and $c_b$ is the Lloyd-Max distortion coefficient at $b$ bits.
\end{theorem}

\begin{proof}
The compressed output decomposes as:
\begin{equation}
    \hat{o}_{\text{quant}} = \hat{a}^\top \hat{V} = a^\top V + \underbrace{a^\top \delta_V}_{E_V} + \underbrace{(\hat{a} - a)^\top \hat{V}}_{E_K}
\end{equation}

\paragraph{Value error ($E_V$).}
The Value quantization errors $\delta_{v_i} = \hat{v}_i - v_i$ are independent across tokens after rotation, with $\mathbb{E}[\|\delta_{v_i}\|^2] = d \cdot c_{b_V} \sigma_V^2$. Therefore:
\begin{equation}
    \mathbb{E}\left[\left\|\sum_i a_i \delta_{v_i}\right\|^2\right] = \sum_i a_i^2 \cdot d \cdot c_{b_V} \sigma_V^2 = d \cdot c_{b_V} \sigma_V^2 \cdot \kappa
\end{equation}

\paragraph{Key error ($E_K$).}
Key quantization perturbs the attention logits: $\hat{s}_i - s_i = q^\top \delta_{k_i} / \sqrt{d}$, with variance $\sigma_s^2 = \sigma_Q^2 c_{b_K} \sigma_K^2$. The softmax amplification is bounded by $\|\partial a / \partial s\|_F \leq 1/T$, giving:
\begin{equation}
    \mathbb{E}\left[\|(\hat{a} - a)^\top \hat{V}\|^2\right] \leq \frac{\sigma_Q^2 c_{b_K} \sigma_K^2}{T^2} \cdot d\sigma_V^2
\end{equation}

Combining both terms yields the stated bound.
\end{proof}

\begin{remark}[K sensitivity is amplified by GQA]
The $E_K$ term scales with $1/T^2$, capturing the softmax amplification of Key errors. In a GQA setting with ratio $r = n_q / n_\text{kv}$, each KV head serves $r$ query heads, so the total K-error contribution to the layer output is amplified by a factor of $r$. For $r = 7$ (Qwen), K-quantization error at 4-bit crosses the catastrophic threshold, while V$=$4-bit remains safe because V errors are linearly averaged (no softmax amplification). The sensitivity ranking K $>$ V is \emph{architecture-dependent}: it becomes more pronounced as GQA ratio increases.
\end{remark}

\subsection{Eviction Error Bound}

\begin{theorem}[Eviction Threshold]
\label{thm:evict}
A token $j$ can be evicted with bounded error when its attention-weighted contribution falls below the quantization noise floor:
\begin{equation}
    a_j < \tau_{\text{evict}} = \alpha \cdot \sqrt{\frac{c_{b_V} \kappa}{d} + \frac{c_{b_K} \sigma_Q^2 \sigma_K^2}{T^2 d}}
\end{equation}
where $\alpha \in (0, 1]$ is a safety margin. For $b_V = 4$, $d = 128$: $\tau_{\text{evict}} \approx 0.017\alpha$.
\end{theorem}

\begin{theorem}[Total Eviction Error]
\label{thm:evict-total}
If all evicted tokens satisfy the threshold criterion (Theorem~\ref{thm:evict}), and $A_{\mathcal{E}} = \sum_{j \in \mathcal{E}} a_j$ is the total evicted attention mass, then:
\begin{equation}
    E_{\text{evict}} \leq A_{\mathcal{E}} \cdot \tau_{\text{evict}} \cdot d \sigma_V^2
\end{equation}
For typical transformer attention distributions, the bottom 50\% of middle tokens carry $<$5\% of total attention mass, making the eviction error negligible compared to quantization error.
\end{theorem}

\subsection{Joint Error Bound}

\begin{theorem}[Joint Error Bound]
\label{thm:joint}
Under the combined quantization and eviction pipeline:
\begin{equation}
    E_{\text{total}} \leq E_{\text{quant}}(\mathcal{S}) + E_{\text{evict}} + 2\sqrt{E_{\text{quant}}(\mathcal{S}) \cdot E_{\text{evict}}}
\end{equation}
When eviction targets tokens with small attention mass ($A_{\mathcal{E}} \ll 1$), the cross-term is negligible and the errors are approximately additive:
\begin{equation}
    E_{\text{total}} \approx E_{\text{quant}} + A_{\mathcal{E}}^2 \cdot d\sigma_V^2
\end{equation}
\end{theorem}

\subsection{Threshold Derivation}
\label{sec:threshold-derivation}

The four-tier thresholds $\tau_1 < \tau_2 < \tau_3$ are derived from the quantization noise floor $\sigma_{\text{floor}}$:

\paragraph{$\tau_1$ (eviction).} From Theorem~\ref{thm:evict}: $\tau_1 = \alpha_1 \cdot \sigma_{\text{floor}} / (\sigma_V \sqrt{d})$.

\paragraph{$\tau_2$ (V degradation).} Token $j$ tolerates V$=$2-bit when the error difference is within the noise floor:
\begin{equation}
    \tau_2 = \frac{\alpha_2}{\sqrt{c_2} - \sqrt{c_4}} \cdot \frac{\sigma_{\text{floor}}}{\sigma_V\sqrt{d}} = \frac{\alpha_2}{0.414} \cdot \frac{\tau_1}{\alpha_1} \approx 2.4 \cdot \frac{\alpha_2}{\alpha_1} \cdot \tau_1
\end{equation}

\paragraph{$\tau_3$ (V precision upgrade).} Same derivation between V$=$4-bit and V$=$6-bit:
\begin{equation}
    \tau_3 = \frac{\alpha_3}{\sqrt{c_4} - \sqrt{c_6}} \cdot \frac{\sigma_{\text{floor}}}{\sigma_V\sqrt{d}} = \frac{\alpha_3}{0.108} \cdot \frac{\tau_1}{\alpha_1} \approx 9.3 \cdot \frac{\alpha_3}{\alpha_1} \cdot \tau_1
\end{equation}

With $\alpha_1 = \alpha_2 = \alpha_3 = 1$: $\tau_1 : \tau_2 : \tau_3 = 1 : 2.4 : 9.3$.

\subsection{The PPL--Task Evaluation Gap}
\label{sec:ppl-gap}

Our L2 error bounds predict higher safe eviction rates than observed under PPL evaluation. The gap arises because PPL $= \exp\left(\frac{1}{N}\sum_i \ell_i\right)$ is a geometric mean of per-token losses---a single token with large cross-entropy increase disproportionately affects the aggregate.

Formally, the cross-entropy increase for token $i$ is:
\begin{equation}
    \Delta\ell_i \approx \frac{\Delta p_i}{p_i}
\end{equation}
For low-confidence predictions ($p_i$ small), even modest perturbations cause large $\Delta\ell_i$. These are precisely the tokens where evicted context matters most, creating a systematic amplification.

\paragraph{Context length scaling.}
The safe eviction rate under PPL follows logarithmic scaling:
\begin{equation}
    f_{\text{safe}}(n) = 0.277 + 0.0405 \cdot \ln(n)
\end{equation}
This is slower than the na\"ive $1 - O(1/n)$ prediction because the number of semantically important middle tokens grows as $O(\log n)$, reflecting hierarchical document structure.

\paragraph{Analogy to video compression.}
The PPL--task gap mirrors the PSNR--VMAF gap in video compression. PSNR penalizes every pixel equally; VMAF weights by perceptual importance. Similarly, PPL penalizes every token prediction equally, while task metrics (NIAH, LongBench) only evaluate task-relevant tokens. This typically allows 2--3$\times$ higher compression under perceptual/task metrics at equivalent quality.

\subsection{Information-Theoretic Analysis}
\label{sec:info-theory}

\paragraph{Shannon rate-distortion bound.}
After Hadamard rotation, each KV entry is approximately $\mathcal{N}(0, \sigma^2 I_d)$. The Shannon rate-distortion function gives the minimum bits per entry:
\begin{equation}
    R(D) = \frac{d}{2} \log_2\left(\frac{\sigma^2}{D}\right)
\end{equation}

For our distortion level ($D = c_6 \sigma^2 = 0.0066\sigma^2$):
\begin{equation}
    R(D) = \frac{128}{2} \times 7.24 = 463 \text{ bits/entry} = 3.62 \text{ bits/dim}
\end{equation}

TurboQuant uses 6 bits/dim, yielding a coding efficiency of $\eta = 3.62/6 = 60.3\%$. The 40\% gap is the well-known ``space-filling loss'' of scalar vs.\ vector quantization~\citep{conway1999}.

\paragraph{Compression ceiling.}
Combining optimal quantization with eviction at $n = 128K$ ($f_{\text{evict}} = 0.75$):
\begin{align}
    \rho_{\text{Shannon}} &= \frac{32}{(1 - 0.75) \times 2 \times 3.01} = 21.3\times \quad \text{(Shannon limit)} \\
    \rho_{\text{practical}} &= \frac{32}{(1 - 0.75) \times 2 \times 5.0} = 12.8\times \quad \text{(at 60\% efficiency)}
\end{align}

Our measured ceiling of $\sim$13$\times$ matches the practical rate-distortion bound, confirming we operate at $\sim$80\% of the theoretical limit for PPL-constrained compression.

\paragraph{Task-metric ceiling.}
Under task-based evaluation, the effective eviction constraint relaxes to $f_{\text{evict}} \leq 0.95$ (only task-relevant tokens need preservation):
\begin{equation}
    \rho_{\text{task}} = \frac{32}{0.05 \times 2 \times 5.0} = 64\times
\end{equation}

The 30$\times$ target lies comfortably between the PPL limit (13$\times$) and the task limit (64$\times$).

\paragraph{Reverse water-filling and four-tier optimality.}
The optimal bit allocation across sources with different importance is given by the reverse water-filling algorithm~\citep{cover2006}:
\begin{equation}
    R_k = \max\left(0, \frac{1}{2}\log_2\frac{\sigma_k^2}{\theta}\right)
\end{equation}
where $\theta$ is the water level determined by the rate budget. In our framework, each token's ``importance'' is its attention weight $a_j$. The four-tier classification is the discrete approximation: tokens below the water level ($a_j \leq \tau_1$) receive zero bits (eviction), while tokens above receive progressively more bits (8, 10, or 12) based on their importance tier.

\begin{table}[h]
\centering
\begin{tabular}{lcc}
\toprule
\textbf{Metric} & \textbf{Compression ceiling} & \textbf{Our efficiency} \\
\midrule
Shannon limit (PPL) & 21$\times$ & --- \\
Practical limit (PPL, 60\% eff.) & 13$\times$ & $\sim$80\% \\
Shannon limit (task) & 64$\times$+ & --- \\
30$\times$ target & --- & Within reach \\
\bottomrule
\end{tabular}
\caption{Information-theoretic compression limits. Our PPL ceiling matches the practical rate-distortion bound, proving the limit is fundamental.}
\label{tab:info-theory}
\end{table}
